\chapter{The scene, meshes, undo and assets}
\label{ch:scene}

\section{The mesh}
\code{Mesh} (\code{src/scene/mesh.h}) keeps Blender's layout \cite{Blender2026}: \code{positions}, \code{face\_offsets}
and \code{corner\_verts}, so a face is a range of corners and n-gons are first-class; per-corner UVs, per-face material
indices, seams, sharp edges, loose edges, custom normals, vertex groups and shape keys are named attributes beside them
\cite{MarschnerShirley2021}. The render cache splits corners into render vertices only where their attributes differ
(\code{build\_render\_vertices}, bit-identical in parallel and serial), builds triangles per face range and MikkTSpace
tangents \cite{Mikkelsen2008Simulation}; a positions-only change follows the kept topology in place
(\code{update\_in\_place}, \adr{0034}), with the deforming tail of a modifier stack re-run on a kept mesh (\adr{0022}).

Meshes are shared through \code{MeshPtr} as Blender's implicit sharing does (\adr{0003}). An undo step is a clone of the
scene that shares unchanged meshes; code that edits a mesh calls \code{mesh\_make\_mutable} first; every change bumps
\code{Mesh::version}; caches key on pointer + version or on \code{RenderMesh::serial}. Holding a \code{Mesh\&} across an
operation that may replace the mesh is a bug reviewers look for.

\section{The operators}
\code{meshops} (\code{src/scene/mesh.cpp}, \code{mesh\_tools.cpp} to \code{mesh\_tools5.cpp}) is the Edit Mode toolbox,
each a pure function on a mesh and selection vectors: extrude (faces, regions by wedge, edges, vertices), inset
(individual and region with a continuous limit), bevel with a superellipse profile and clamp overlap, bridge with a
Hermite tube, loop cut, knife, connect, dissolve (edges, limited, degenerate), collapse, merge (at centre, by distance,
selected-to-unselected), fill and smart fill, poke, triangulate and tris-to-quads, Catmull--Clark
\cite{Catmull1978Recursively}, Taubin smoothing \cite{Taubin1995Signal}, spin, slice, follow, shell, draft, radial
array, fillet, grid, pipe, array, taper, recess, and the SketchUp family: push/pull with its limits, flush steps and
holes (\adr{0004}), push through, attach face, imprint across faces, repair rings. \code{src/scene/mesh\_internal.h}
holds the shared topology helpers (face builder, edge-faces table, the \code{annulus} ring builder).

Robustness is a requirement (Q-01), not a hope: the stress suite (\code{stress/stress\_main.cpp}) runs the operators on
odd meshes (degenerate faces, non-manifold edges, open surfaces, extreme scales, NaN) and counts problems and faces on
top of each other (\code{overlapping\_faces}, by Sutherland--Hodgman clipping \cite{Sutherland1974}); the counts may not
rise without an explanation in the CHANGELOG.

\section{Objects and components}
As in Unity, everything in a scene is a \code{GameObject} with a \code{Transform}; what it does comes from components:
\code{MeshFilter}, \code{MeshRenderer}, \code{Light}, \code{Camera}, \code{Rigidbody}, \code{Animation},
\code{Armature}, \code{Bone}, modifiers (\code{SubdivisionSurface}, \code{Mirror}, \code{Array}, \code{Boolean}, \ldots,
\adr{0022}), \code{CameraFilters} (\adr{0008}), \code{ProbeVolume}, \code{ReflectionProbe}, \code{SpringBone}
(\adr{0049}), \code{Collider}, \code{LODGroup}, \code{AvatarDescriptor} and VRChat's others (\adr{0100}), and
\code{EngineComponent} for what another engine's file held that Blendity does not interpret (\adr{0101}). A component
says what the views read of it (\code{render\_role}), so the scene hash takes one virtual call instead of a
\code{dynamic\_cast} per type \cite{Gregory2018vol2}.

\section{The Reflector}
A component describes its fields once, in \code{reflect(Reflector\&)}: name, type, range, help, groups, typed object
and asset references. That one description draws the Inspector (with essentials first, multi-object editing and
expressions in fields), writes and reads the scene file (tolerant of unknown fields, so files move between versions),
hashes the component for change detection, answers the console's \code{set Component.Field value}, keys any field from
the Animation window (\rnd{335}), and carries fields over Live Link and into Unity's write-back. It is Blender's RNA
\cite{Blender2026} and Unity's \code{SerializedProperty} in one. Lists of effects or layers (a camera's filter stack, a
texture's layers) are reflected through a prefix per element, so files, undo and hashing work through the same path.

\section{Scene IO and assets}
Scenes are a text format (\code{src/scene/scene\_io.cpp}): one line per object and field, materials and textures by
asset path with values kept as a fallback, lightmaps and probes beside the scene (\adr{0010}). Material assets
(\code{.mat}) and procedural textures (\code{.ptex}, \adr{0017}) are shared instances written back after each finished
change; after loading, the file wins. The Project window moves files on disk and retargets every reference (the asset
library, texture paths, scene files, clips); deletes go to the OS trash; moves and deletes are recorded so an undo that
brings back an old path follows the file. Imported models (OBJ, FBX through ufbx, glTF, STL, PLY, USD, PMX, VRM, BVH,
Collada, \code{.blend} through Blender) arrive with materials as \code{.mat} assets, textures found or extracted, skins
and clips (\adr{0047}); exports share one preparation step (world transform, axes, triangulation, split normals,
baked textures) and the writers for each format \cite{Khronos2021Gltf, FBXBlender2026, USD2026}.

\section{Undo}
\code{UndoState} snapshots the scene graph after each finished interaction (\code{commit\_undo}, with a label the
History window lists); meshes are shared and copied on first write; a sculpt or paint stroke is a lighter delta step of
the vertices or the 64$\times$64 texture tiles it touched (\adr{0074}, \adr{0081}); the Draw tab's paper and the reference
board keep histories of their own. Vertex groups and custom normals are mended by position at each commit so a
topology change does not lose them (\adr{0024}). Undo covers the scene, materials, textures and clips as one step per
operation (Q-05), and the budget of kept steps is bytes, oldest out (\adr{0026}).
